\documentclass[12pt, titlepage]{article}

\usepackage{amsmath, mathtools}

\usepackage[round]{natbib}
\usepackage{amsfonts}
\usepackage{amssymb}
\usepackage{graphicx}
\usepackage{colortbl}
\usepackage{xr}
\usepackage{hyperref}
\usepackage{longtable}
\usepackage{xfrac}
\usepackage{tabularx}
\usepackage{float}
\usepackage{siunitx}
\usepackage{booktabs}
\usepackage{multirow}
\usepackage[section]{placeins}
\usepackage{caption}
\usepackage{fullpage}

\hypersetup{
bookmarks=true,     % show bookmarks bar?
colorlinks=true,       % false: boxed links; true: colored links
linkcolor=red,          % color of internal links (change box color with linkbordercolor)
citecolor=blue,      % color of links to bibliography
filecolor=magenta,  % color of file links
urlcolor=cyan          % color of external links
}

\usepackage{array}

\externaldocument{../../SRS/SRS}

\input{../../Comments}
%% Common Parts

\newcommand{\progname}{Cut Edge Weight: Data Separability Measurement} % PUT YOUR PROGRAM NAME HERE
\newcommand{\authname}{Lesley Wheat} % AUTHOR NAMES                  

\usepackage{hyperref}
    \hypersetup{colorlinks=true, linkcolor=blue, citecolor=blue, filecolor=blue,
                urlcolor=blue, unicode=false}
    \urlstyle{same}
                                


\begin{document}

\title{Module Interface Specification for \progname{}}

\author{\authname}

\date{\today}

\maketitle

\pagenumbering{roman}

\section{Revision History}

\begin{tabularx}{\textwidth}{p{3cm}p{2cm}X}
\toprule {\bf Date} & {\bf Version} & {\bf Notes}\\
\midrule
2022-03-10 & 1.0 & Creation \\
\bottomrule
\end{tabularx}

~\newpage

\section{Symbols, Abbreviations and Acronyms}

See SRS Documentation at \url{https://github.com/LesleyWheat/CAS741_CEWS/blob/main/docs/SRS/SRS.pdf} .

\mycomment{Also add any additional symbols, abbreviations or acronyms}

\newpage

\tableofcontents

\newpage

\pagenumbering{arabic}

\section{Introduction}

The following document details the Module Interface Specifications for
\mycomment{Fill in your project name and description}

Complementary documents include the System Requirement Specifications
and Module Guide.  The full documentation and implementation can be
found at \url{https://github.com/LesleyWheat/CAS741_CEWS/blob/main/docs/SRS/SRS.pdf}.

\section{Notation}

\mycomment{You should describe your notation.  You can use what is below as
  a starting point.}

The structure of the MIS for modules comes from \citet{HoffmanAndStrooper1995},
with the addition that template modules have been adapted from
\cite{GhezziEtAl2003}.  The mathematical notation comes from Chapter 3 of
\citet{HoffmanAndStrooper1995}.  For instance, the symbol := is used for a
multiple assignment statement and conditional rules follow the form $(c_1
\Rightarrow r_1 | c_2 \Rightarrow r_2 | ... | c_n \Rightarrow r_n )$.

The following table summarizes the primitive data types used by \progname. 

\begin{center}
\renewcommand{\arraystretch}{1.2}
\noindent 
\begin{tabular}{l l p{7.5cm}} 
\toprule 
\textbf{Data Type} & \textbf{Notation} & \textbf{Description}\\ 
\midrule
character & char & a single symbol or digit\\
integer & $\mathbb{Z}$ & a number without a fractional component in (-$\infty$, $\infty$) \\
natural number & $\mathbb{N}$ & a number without a fractional component in [1, $\infty$) \\
real & $\mathbb{R}$ & any number in (-$\infty$, $\infty$)\\
\bottomrule
\end{tabular} 
\end{center}

\noindent
The specification of \progname \ uses some derived data types: sequences, strings, and
tuples. Sequences are lists filled with elements of the same data type. Strings
are sequences of characters. Tuples contain a list of values, potentially of
different types. In addition, \progname \ uses functions, which
are defined by the data types of their inputs and outputs. Local functions are
described by giving their type signature followed by their specification.

\newpage
\section{Module Decomposition}

The following table is taken directly from the Module Guide document for this project.

\begin{table}[h!]
\centering
\begin{tabular}{p{0.3\textwidth} p{0.6\textwidth}}
\toprule
\textbf{Level 1} & \textbf{Level 2}\\
\midrule

{Hardware-Hiding Module} & ~ \\
\midrule

\multirow{3}{0.3\textwidth}{Behaviour-Hiding Module} 
& Control Module\\
& Input Module\\
& Output Module\\

\midrule

\multirow{2}{0.3\textwidth}{Software Decision Module} 
& Calculation Module\\
& Graph Module\\
\bottomrule

\end{tabular}
\caption{Module Hierarchy}
\label{TblMH}
\end{table}

\mycomment{
\newpage
~\newpage

\section{MIS of \mycomment{Module Name}} \label{Module} \mycomment{Use labels for
  cross-referencing}

\mycomment{You can reference SRS labels, such as R\ref{R_Inputs}.}

\mycomment{It is also possible to use \LaTeX for hypperlinks to external documents.}

\subsection{Module}

\mycomment{Short name for the module}

\subsection{Uses}


\subsection{Syntax}

\subsubsection{Exported Constants}

\subsubsection{Exported Access Programs}

\begin{center}
\begin{tabular}{p{2cm} p{4cm} p{4cm} p{2cm}}
\hline
\textbf{Name} & \textbf{In} & \textbf{Out} & \textbf{Exceptions} \\
\hline
\mycomment{accessProg} & - & - & - \\
\hline
\end{tabular}
\end{center}

\subsection{Semantics}

\subsubsection{State Variables}

\mycomment{Not all modules will have state variables.  State variables give the module
  a memory.}

\subsubsection{Environment Variables}

\mycomment{This section is not necessary for all modules.  Its purpose is to capture
  when the module has external interaction with the environment, such as for a
  device driver, screen interface, keyboard, file, etc.}

\subsubsection{Assumptions}

\mycomment{Try to minimize assumptions and anticipate programmer errors via
  exceptions, but for practical purposes assumptions are sometimes appropriate.}

\subsubsection{Access Routine Semantics}

\noindent \mycomment{accessProg}():
\begin{itemize}
\item transition: \mycomment{if appropriate} 
\item output: \mycomment{if appropriate} 
\item exception: \mycomment{if appropriate} 
\end{itemize}

\mycomment{A module without environment variables or state variables is unlikely to
  have a state transition.  In this case a state transition can only occur if
  the module is changing the state of another module.}

\mycomment{Modules rarely have both a transition and an output.  In most cases you
  will have one or the other.}

\subsubsection{Local Functions}

\mycomment{As appropriate} \mycomment{These functions are for the purpose of specification.
  They are not necessarily something that is going to be implemented
  explicitly.  Even if they are implemented, they are not exported; they only
  have local scope.}
}

% MM
\newpage
  \section{MIS of Control Module} \label{M_CO}

\subsection{Control Module}

\subsection{Uses}

Input Module, Calculation Module, Output Module

\subsection{Syntax}

\subsubsection{Exported Constants}
None.

\subsubsection{Exported Access Programs}

\begin{center}
\begin{tabular}{p{2cm} p{4cm} p{4cm} p{2cm}}
\hline
\textbf{Name} & \textbf{In} & \textbf{Out} & \textbf{Exceptions} \\
\hline
Main & data, labels & $J_N$ & - \\
\hline
\end{tabular}
\end{center}

\subsection{Semantics}

\subsubsection{State Variables}
None.

\subsubsection{Environment Variables}
None.

\subsubsection{Assumptions}

\subsubsection{Access Routine Semantics}

\noindent \mycomment{accessProg}():
\begin{itemize}
\item transition: \mycomment{if appropriate} 
\item output: \mycomment{if appropriate} 
\item exception: \mycomment{if appropriate} 
\end{itemize}


\subsubsection{Local Functions}
None.

% MM_
\newpage
  \section{MIS of Input Module} \label{M_I}

\subsection{Input Module}

\subsection{Uses}
None.

\subsection{Syntax}

\subsubsection{Exported Constants}
None.

\subsubsection{Exported Access Programs}

\begin{center}
\begin{tabular}{p{2cm} p{4cm} p{4cm} p{2cm}}
\hline
\textbf{Name} & \textbf{In} & \textbf{Out} & \textbf{Exceptions} \\
\hline
processInput & data, labels & - & Incorrect type, data size mismatch \\
\hline
\end{tabular}
\end{center}

\subsection{Semantics}

\subsubsection{State Variables}
None.

\subsubsection{Environment Variables}
None.

\subsubsection{Assumptions}

\subsubsection{Access Routine Semantics}

\noindent \mycomment{accessProg}():
\begin{itemize}
\item transition: \mycomment{if appropriate} 
\item output: \mycomment{if appropriate} 
\item exception: \mycomment{if appropriate} 
\end{itemize}


\subsubsection{Local Functions}


% MM
\newpage
  \section{MIS of Output Module} \label{M_O}

\subsection{Output Module}

\subsection{Uses}
None.

\subsection{Syntax}
\subsubsection{Exported Constants}
None.

\subsubsection{Exported Access Programs}

\begin{center}
\begin{tabular}{p{2cm} p{4cm} p{4cm} p{2cm}}
\hline
\textbf{Name} & \textbf{In} & \textbf{Out} & \textbf{Exceptions} \\
\hline
checkOutput & $J_N$ & - & resultOutOfBounds \\
\hline
\end{tabular}
\end{center}

\subsection{Semantics}
\subsubsection{State Variables}
None.

\subsubsection{Environment Variables}
None.

\subsubsection{Assumptions}

\subsubsection{Access Routine Semantics}

\noindent \mycomment{accessProg}():
\begin{itemize}
\item transition: \mycomment{if appropriate} 
\item output: \mycomment{if appropriate} 
\item exception: \mycomment{if appropriate} 
\end{itemize}


\subsubsection{Local Functions}
None.

% MM
\newpage
\section{MIS of Calculation Module} \label{M_CA}

\subsection{Calculation Module}

\subsection{Uses}
Graph Module

\subsection{Syntax}
\subsubsection{Exported Constants}
None.

\subsubsection{Exported Access Programs}

\begin{center}
\begin{tabular}{p{2cm} p{4cm} p{4cm} p{2cm}}
\hline
\textbf{Name} & \textbf{In} & \textbf{Out} & \textbf{Exceptions} \\
\hline
calcCEWS & data, labels & $J_N$ & - \\
\hline
\end{tabular}
\end{center}

\subsection{Semantics}
\subsubsection{State Variables}
None.

\subsubsection{Environment Variables}
None.

\subsubsection{Assumptions}


\subsubsection{Access Routine Semantics}

\noindent \mycomment{accessProg}():
\begin{itemize}
\item transition: \mycomment{if appropriate} 
\item output: \mycomment{if appropriate} 
\item exception: \mycomment{if appropriate} 
\end{itemize}


\subsubsection{Local Functions}
None.

% MM
\newpage
\section{MIS of Graph Module} \label{M_G}

\subsection{Graph Module}

\subsection{Uses}
None.

\subsection{Syntax}
\subsubsection{Exported Constants}
None.

\subsubsection{Exported Access Programs}

\begin{center}
\begin{tabular}{p{2cm} p{4cm} p{4cm} p{2cm}}
\hline
\textbf{Name} & \textbf{In} & \textbf{Out} & \textbf{Exceptions} \\
\hline
makeGraph & data & graphObject & - \\
\hline
\end{tabular}
\end{center}

\subsection{Semantics}
\subsubsection{State Variables}
None.

\subsubsection{Environment Variables}
None.

\subsubsection{Assumptions}

\subsubsection{Access Routine Semantics}

\noindent \mycomment{accessProg}():
\begin{itemize}
\item transition: \mycomment{if appropriate} 
\item output: \mycomment{if appropriate} 
\item exception: \mycomment{if appropriate} 
\end{itemize}


\subsubsection{Local Functions}
None.

\newpage

\bibliographystyle {plainnat}
\bibliography {../../../refs/References}

\newpage

\section{Appendix} \label{Appendix}

\mycomment{Extra information if required}

\end{document}